% Weekly-Template.tex — IEEEtran 双栏中文周报模板（weekly-report）
% 编译方式：XeLaTeX（本地 compile.sh，或 Overleaf 选 XeLaTeX）。
% build_pdf.py 只读取本文件 \begin{document} 之前的导言区，正文由代码按 facts.json 生成。
\documentclass[journal]{IEEEtran}
\usepackage[usenames]{color}
\usepackage{graphicx}
\usepackage{amsmath}
\usepackage{amssymb}
\usepackage{multirow}
\usepackage{array}
\usepackage{url}
\usepackage{xcolor}
\usepackage{hhline}
\usepackage{booktabs}
\usepackage[letterpaper]{geometry}
\geometry{verbose,tmargin=0.62in,bmargin=0.9in,lmargin=0.65in,rmargin=0.65in}
\setlength{\headheight}{18pt}
\setlength{\headsep}{10pt}

% 中文支持（需 XeLaTeX；compile.sh 会把字体降级到本机实际存在的中文字体）
\usepackage{xeCJK}
\setCJKmainfont{Noto Serif CJK SC}
\setCJKsansfont{Noto Sans CJK SC}
\setCJKmonofont{Noto Sans Mono CJK SC}

\renewcommand{\abstractname}{\textbf{结论 / SUMMARY}}
\usepackage{fancyhdr}
\pagestyle{fancy}
\fancyhf{}
\lhead{\small WEEKLY REPORT / 工作周报}
\rhead{\small \thepage}

\begin{document}
% ---- 以下为占位示例，实际正文由 build_pdf.py 生成 ----
\title{这里是周报标题}
\author{作者 / 周期}
\maketitle
\begin{abstract}
这里是一句话结论。
\end{abstract}
\section{已完成 / DONE}
这里是条目。
\begin{figure}[t]
\centering
\includegraphics[width=0.92\linewidth]{./fig1.png}
\caption{图的说明（必须写明数据来源）。}
\end{figure}
\end{document}
